% ===================== ANEXOS =====================
\appendix
\section{[Titulo del Anexo Tecnico]}
\label{appendix:calculations}

[Incluir en este anexo los calculos detallados, desarrollos matematicos,
balances de masa o energia, y modelos fisicos que sustenten los resultados
presentados en el cuerpo del informe. Cada ecuacion debe estar numerada,
con las variables definidas inmediatamente despues.]

\subsection{[Nombre del Modelo o Calculo]}

[Describir el objetivo del calculo, los supuestos del modelo y su rango
de validez. Referenciar la fuente bibliografica del modelo.]

\subsubsection{Ecuaciones fundamentales}

[Presentar las ecuaciones fundamentales del modelo. Ejemplo de balance termico:]

La resistencia termica por conveccion interna se expresa como:
\begin{equation}
    R_{\mathrm{int}} = \frac{1}{h_{\mathrm{int}} \, \pi D_{\mathrm{int}} L},
    \label{eq:r_int}
\end{equation}
donde $h_{\mathrm{int}}$ es el coeficiente de transferencia de calor por conveccion
forzada, $D_{\mathrm{int}}$ el diametro interno y $L$ la longitud de la tuberia.

La resistencia por conduccion en la pared metalica:
\begin{equation}
    R_{\mathrm{pared}} = \frac{\ln(D_{\mathrm{ext}}/D_{\mathrm{int}})}{2\pi k_{\mathrm{acero}} L},
    \label{eq:r_pared}
\end{equation}
con $k_{\mathrm{acero}}$ la conductividad termica del material.

El flujo de calor perdido hacia el ambiente:
\begin{equation}
    \dot{Q}_{\mathrm{perdida}} = \frac{T_{\mathrm{fluido}} - T_{\infty}}{R_{\mathrm{total}}}.
    \label{eq:q_perdida}
\end{equation}

La caida de temperatura del fluido:
\begin{equation}
    \Delta T = \frac{\dot{Q}_{\mathrm{perdida}}}{\dot{m} \, c_{p}},
    \label{eq:delta_t}
\end{equation}
donde $\dot{m}$ es el flujo masico y $c_{p}$ la capacidad calorifica especifica.

\subsubsection{Parametros de entrada y supuestos}

[Incluir tabla con los parametros fisicos, geometricos y de operacion
utilizados en el calculo, con sus fuentes.]

\begin{table}[htbp]
    \centering
    \caption{[Leyenda descriptiva de parametros del modelo.]}
    \label{tab:param_modelo}
    \begin{tabular}{@{}llcl@{}}
        \toprule
        \textbf{Parametro} & \textbf{Simbolo} & \textbf{Valor} & \textbf{Unidad} \\
        \midrule
        <Parametro 1> & <$S_1$> & <Valor> & <Unidad> \\
        <Parametro 2> & <$S_2$> & <Valor> & <Unidad> \\
        <Parametro 3> & <$S_3$> & <Valor> & <Unidad> \\
        \bottomrule
    \end{tabular}
\end{table}

\section{[Titulo del Segundo Anexo]}
\label{appendix:drawings}

[Relacionar planos, diagramas, isometricos o documentacion tecnica consultada
como base para el estudio. Incluir tabla de relacion de documentos.]

\begin{table}[htbp]
    \centering
    \caption{[Leyenda descriptiva de planos tecnicos consultados.]}
    \label{tab:planos}
    \begin{tabular}{@{}lll@{}}
        \toprule
        \textbf{Documento} & \textbf{Descripcion} & \textbf{Ubicacion} \\
        \midrule
        <Documento 1> & <Descripcion> & <Ruta / Carpeta> \\
        <Documento 2> & <Descripcion> & <Ruta / Carpeta> \\
        \bottomrule
    \end{tabular}
\end{table}

\section{[Titulo del Tercer Anexo]}
\label{appendix:specs}

[Incluir especificaciones tecnicas de equipos relevantes para el estudio.
Una subseccion por equipo.]

\subsection{[Equipo 1]}

\begin{table}[htbp]
    \centering
    \caption{[Leyenda descriptiva de especificaciones del equipo.]}
    \label{tab:spec_equipo1}
    \begin{tabular}{@{}lll@{}}
        \toprule
        \textbf{Parametro} & \textbf{Valor} & \textbf{Unidad} \\
        \midrule
        <Fabricante / Modelo> & <Valor> & <---> \\
        <Parametro 1> & <Valor> & <Unidad> \\
        <Parametro 2> & <Valor> & <Unidad> \\
        \bottomrule
    \end{tabular}
\end{table}
